\chapter{Equivalence Checking and Formal Verification}
\label{ch:proof}

\glance{This chapter surveys the verification engines that live under \pkg{src/proof}: combinational equivalence checking (CEC), functional reduction (FRAIG), structural choice computation for mapping, and the sequential model-checking family (PDR, signal correspondence, interpolation, abstraction). All of them ultimately reduce a question about circuit behaviour to a sequence of SAT queries, and all of them report their verdict through a shared counter-example structure, \id{Abc_Cex_t}.}

%===============================================================
\section{The Miter: Reducing Equivalence to SAT}
%===============================================================

The unifying idea behind equivalence checking is the \emph{miter}. Given two
circuits with matching inputs and outputs, ABC builds a new circuit that
XOR-s the corresponding outputs of the two designs and OR-s the XORs into a
single output. The two circuits are equivalent \emph{iff} that output can
never be~1: a SAT solver that proves the output unsatisfiable establishes
equivalence, while a satisfying assignment is a concrete input pattern on
which the circuits differ (a counter-example).

The classic network-level constructor is \id{Abc_NtkMiter} in
\file{src/base/abci/abcMiter.c}:

\begin{lstlisting}[style=abccode]
// src/base/abci/abcMiter.c
Abc_Ntk_t * Abc_NtkMiter( Abc_Ntk_t * pNtk1, Abc_Ntk_t * pNtk2,
                          int fComb, int nPartSize,
                          int fImplic, int fMulti );
\end{lstlisting}

Callers such as \file{src/base/abci/abcDar.c} build the miter and then ask
whether it is constant~0:

\begin{lstlisting}[style=abccode]
// src/base/abci/abcDar.c (Abc_NtkDarCec path)
pMiter   = Abc_NtkMiter( pNtk1, pNtk2, 0, 0, 0, 0 );
RetValue = Abc_NtkMiterIsConstant( pMiter );
\end{lstlisting}

Modern GIA-based engines fold the miter directly into the solver: functions
such as \id{Cec_ManVerifyTwo} in \file{src/proof/cec/cecCec.c} build an
internal miter from two \id{Gia_Man_t} objects and call \id{Cec_ManVerify}.

%===============================================================
\section{Combinational Equivalence Checking (\pkg{src/proof/cec})}
%===============================================================

The \pkg{src/proof/cec} package performs combinational equivalence checking on
And-Inverter Graphs. Rather than throwing the whole miter at a monolithic SAT
call, it interleaves two complementary techniques, coordinated from
\file{src/proof/cec/cecCore.c}:

\begin{itemize}
  \item \textbf{Random/guided simulation} groups nodes into candidate
        equivalence classes (nodes that agree on all simulated patterns).
        A satisfying pattern found by SAT is fed back as new simulation
        stimulus, refining the classes.
  \item \textbf{SAT sweeping} proves or disproves the candidate equivalences
        node by node, from inputs to outputs, merging nodes that are proven
        equal so the graph shrinks as the sweep proceeds.
\end{itemize}

Key entry points in \file{src/proof/cec/cecCore.c}:

\begin{lstlisting}[style=abccode]
// src/proof/cec/cecCore.c
Gia_Man_t * Cec_ManSatSolving ( Gia_Man_t * pAig, Cec_ParSat_t * pPars,
                                int f0Proved );
int         Cec_ManSimulationOne( Gia_Man_t * pAig, Cec_ParSim_t * pPars );
Gia_Man_t * Cec_ManSatSweeping ( Gia_Man_t * pAig, Cec_ParFra_t * pPars,
                                 int fSilent );
\end{lstlisting}

Parameter defaults are set by helpers such as \id{Cec_ManSatSetDefaultParams}
(conflict limit \id{nBTLimit}${=}100$, \id{nSatVarMax}${=}2000$) and
\id{Cec_ManFraSetDefaultParams} (up to \id{nItersMax}${=}10$ SAT-sweeping
iterations). The user-facing verdict routine is \id{Cec_ManVerify} in
\file{src/proof/cec/cecCec.c}.

\begin{table}[t]
\centering
\begin{tabularx}{\linewidth}{@{}lX@{}}
\toprule
\textbf{Command} & \textbf{Purpose} \\
\midrule
\cmd{cec}   & Combinational equivalence check of two logic networks (network side). \\
\cmd{\&cec} & Same check on the GIA (\id{Gia_Man_t}) representation of the \cmd{\&}-space. \\
\cmd{dsat}  & SAT-solve the current network's miter/output directly. \\
\bottomrule
\end{tabularx}
\caption{Combinational equivalence-checking commands.}
\label{tab:cec-cmds}
\end{table}

%===============================================================
\section{FRAIG: Functional Reduction (\pkg{src/proof/fra}, \pkg{src/proof/fraig})}
%===============================================================

A FRAIG (Functionally Reduced AIG) is an AIG in which no two nodes implement
the same Boolean function. Ordinary structural hashing only merges nodes that
are \emph{syntactically} identical; FRAIGing goes further by merging nodes that
are \emph{functionally} identical even when their structure differs. The
procedure combines the same two ingredients as CEC:

\begin{enumerate}
  \item Simulate to find \emph{candidate} equivalences (nodes agreeing on all
        simulation vectors).
  \item For each candidate, call SAT to \emph{prove} the two nodes equal; if
        proven, merge them; if the SAT call returns a distinguishing pattern,
        use it to refine the simulation classes.
\end{enumerate}

The central manager is \id{Fra_Man_t} (\file{src/proof/fra/fra.h}), which holds
the starting manager \id{pManAig}, the resulting \id{pManFraig}, the candidate
class structure \id{Fra_Cla_t}, the simulation manager \id{Fra_Sml_t}, and the
\id{sat_solver}. The core drivers are declared alongside it:

\begin{lstlisting}[style=abccode]
// src/proof/fra/fra.h
extern void        Fra_FraigSweep  ( Fra_Man_t * pManAig );
extern Aig_Man_t * Fra_FraigPerform( Aig_Man_t * pManAig, Fra_Par_t * pPars );
\end{lstlisting}

The older \pkg{src/proof/fraig} package provides the original network-level
FRAIG engine (files \file{fraigMan.c}, \file{fraigSat.c}, \file{fraigTable.c}),
while \pkg{src/proof/fra} is the AIG-manager-based reimplementation.

\begin{table}[t]
\centering
\begin{tabularx}{\linewidth}{@{}llX@{}}
\toprule
\textbf{Command} & \textbf{Backend} & \textbf{Notes} \\
\midrule
\cmd{fraig}  & \id{Abc_NtkDarFraig} (\file{abcDar.c}) & FRAIG the current network. \\
\cmd{ifraig} & incremental FRAIG & iterative/interpolating variant. \\
\cmd{\&fraig}& GIA engine & FRAIG in the \cmd{\&}-space. \\
\bottomrule
\end{tabularx}
\caption{FRAIG commands and their backends.}
\label{tab:fraig-cmds}
\end{table}

%===============================================================
\section{Structural Choices for Mapping (\pkg{src/proof/dch})}
%===============================================================

Technology mapping (Chapter~\ref{ch:map}) produces better results when the
mapper can pick among several \emph{functionally equivalent} AIG structures for
the same logic. The \pkg{src/proof/dch} package (``don't-care''-based choicing)
computes these \emph{choices}: it takes one or more AIG snapshots, proves which
nodes are functionally equivalent using simulation plus SAT (the same
FRAIG-style loop), and records the alternatives as a choice AIG rather than
collapsing them.

The core routine is \id{Dch_ComputeChoices} in
\file{src/proof/dch/dchCore.c}:

\begin{lstlisting}[style=abccode]
// src/proof/dch/dchCore.c
Aig_Man_t * Dch_ComputeChoices( Aig_Man_t * pAig, Dch_Pars_t * pPars )
{
    Dch_Man_t * p = Dch_ManCreate( pAig, pPars );
    p->ppClasses  = Dch_CreateCandEquivClasses( pAig, pPars->nWords,
                                                pPars->fVerbose );
    Dch_ManSweep( p );                 // prove candidates by SAT
    ...
    pResult = Dch_DeriveChoiceAig( pAig, pPars->fSkipRedSupp );
    return pResult;
}
\end{lstlisting}

Defaults from \id{Dch_ManSetDefaultParams} include \id{nWords}${=}8$ simulation
words and \id{nBTLimit}${=}1000$ conflicts. The command \cmd{dch} produces a
choice AIG that a subsequent \cmd{map}/\cmd{if} call can consume; this is the
mechanism behind the \id{fraig\_store}/\id{fraig\_restore} choice scripts that
accumulate several optimized snapshots before mapping.

%===============================================================
\section{Sequential Verification}
%===============================================================

Sequential (multi-cycle) properties require reasoning across time frames. ABC
ships several complementary engines under \pkg{src/proof}, summarized in
Table~\ref{tab:seq-engines}.

\subsection{Property-Directed Reachability (\pkg{src/proof/pdr})}

PDR (also known as IC3, after Aaron Bradley; ABC's implementation follows
Niklas Een) proves safety properties by incrementally building an inductive
sequence of clause sets, one per time frame, without unrolling the whole
transition relation. The entry point is \id{Pdr_ManSolve} in
\file{src/proof/pdr/pdrCore.c}, which drives \id{Pdr_ManSolveInt}. On failure
it attaches a counter-example (via \id{pAig->pSeqModel} or the multi-output
\id{vCexes}); on success it can dump an inductive invariant. The command is
\cmd{pdr}.

\subsection{Signal Correspondence (\pkg{src/proof/ssw})}

Sequential SAT sweeping generalizes FRAIG across time frames: it finds pairs of
flops/nodes that are equivalent under induction and merges them, shrinking the
sequential circuit. The core is \id{Ssw_SignalCorrespondence} in
\file{src/proof/ssw/sswCore.c} (with \id{Ssw_LatchCorrespondence} for the
latch-only variant). These back the \cmd{scorr} and \cmd{lcorr} commands.

\subsection{Interpolation (\pkg{src/proof/int}, \pkg{src/proof/int2})}

Craig interpolation over-approximates reachable states from the proof of a
bounded check, iterating until an inductive invariant is found or a
counter-example appears. The classic engine is
\id{Inter_ManPerformInterpolation} in \file{src/proof/int/intCore.c}; the newer
rewrite is \id{Int2_ManPerformInterpolation} in \file{src/proof/int2/int2Core.c}.

\subsection{Abstraction (\pkg{src/proof/abs})}

When the full model is too large, abstraction hides most of the logic and
refines it only as counter-examples demand. ABC provides gate-level
abstraction, \id{Gia_ManPerformGla} in \file{src/proof/abs/absGla.c}
(command \cmd{gla}), and variable time-frame abstraction, \id{Gia_VtaPerform}
in \file{src/proof/abs/absVta.c} (command \cmd{vta}).

Plain bounded model checking (BMC) is handled by the SAT layer under
\pkg{src/sat} (see Chapter~\ref{ch:sat}).

\begin{table}[t]
\centering
\begin{tabularx}{\linewidth}{@{}llX l@{}}
\toprule
\textbf{Task} & \textbf{Command} & \textbf{Package} & \textbf{Top function} \\
\midrule
Combinational EC (network) & \cmd{cec}    & \pkg{src/proof/cec} & \id{Cec_ManVerify} \\
Combinational EC (GIA)      & \cmd{\&cec}  & \pkg{src/proof/cec} & \id{Cec_ManSatSweeping} \\
Functional reduction        & \cmd{fraig}  & \pkg{src/proof/fra} & \id{Fra_FraigPerform} \\
Structural choices          & \cmd{dch}    & \pkg{src/proof/dch} & \id{Dch_ComputeChoices} \\
Property-directed reach.     & \cmd{pdr}    & \pkg{src/proof/pdr} & \id{Pdr_ManSolve} \\
Signal correspondence        & \cmd{scorr}  & \pkg{src/proof/ssw} & \id{Ssw_SignalCorrespondence} \\
Interpolation                & \cmd{int}    & \pkg{src/proof/int} & \id{Inter_ManPerformInterpolation} \\
Gate-level abstraction       & \cmd{gla}    & \pkg{src/proof/abs} & \id{Gia_ManPerformGla} \\
\bottomrule
\end{tabularx}
\caption{Verification engines: task, command, package, and top-level function.}
\label{tab:seq-engines}
\end{table}

%===============================================================
\section{Counter-Examples and Verdicts}
%===============================================================

Every disproof produces a concrete stimulus stored in the shared
\id{Abc_Cex_t} structure (\file{src/misc/util/utilCex.h}):

\begin{lstlisting}[style=abccode]
// src/misc/util/utilCex.h
struct Abc_Cex_t_
{
    int      iPo;      // failing PO (zero-based)
    int      iFrame;   // failing time-frame (zero-based)
    int      nRegs;    // number of registers in the miter
    int      nPis;     // number of primary inputs
    int      nBits;    // number of bits of data
    unsigned pData[0]; // bits: nRegs + (iFrame+1)*nPis
};
\end{lstlisting}

For combinational checks \id{iFrame} is~0; for sequential checks it records the
cycle at which the property fails, and \id{pData} holds the input (and initial
register) assignment reconstructed cycle by cycle. The active counter-example
is cached on the global frame as \id{pCex} (\id{Abc_Frame_t_},
\file{src/base/main/mainInt.h}), so downstream commands can inspect or simulate it.

Engines report a three-valued verdict, conventionally:

\begin{table}[h]
\centering
\begin{tabularx}{\linewidth}{@{}clX@{}}
\toprule
\textbf{Value} & \textbf{Meaning} & \textbf{Interpretation} \\
\midrule
$1$  & proved     & property holds / circuits equivalent (UNSAT). \\
$0$  & disproved  & counter-example found and stored in \id{pCex} (SAT). \\
$-1$ & undecided  & resource limit hit before a conclusion. \\
\bottomrule
\end{tabularx}
\caption{Verification return codes.}
\label{tab:verdicts}
\end{table}

PDR follows exactly this convention in \id{Pdr_ManSolve}, mapping internal
``unknown'' ($-2$) entries to undecided ($-1$) before returning.

%===============================================================
\section{Worked Example: An Equivalence Sanity Check}
%===============================================================

The following REPL session reads a design, snapshots it, optimizes the copy
with \cmd{resyn2}, and confirms with \cmd{cec} that optimization preserved
functionality:

\begin{lstlisting}[style=abcshell]
abc 01> read design.blif
abc 02> strash
abc 03> write_aiger a.aig
abc 04> resyn2
abc 05> write_aiger b.aig
abc 06> cec a.aig b.aig
Networks are equivalent.
\end{lstlisting}

If the two files differed, \cmd{cec} would instead print that the networks are
NOT equivalent and store the distinguishing pattern in the frame's \id{pCex}.
